% ============================================================================
% T24 - Chương 2 (bản hoàn chỉnh)
% Đề tài ĐT05 - Hệ thống Quản lý Hàng tồn kho & Doanh thu Siêu thị
% Người viết: TV2
%
% File này chỉ chứa NỘI DUNG chương (\chapter ...), \input vào báo cáo chính.
% Cần các gói sau trong preamble của mẫu báo cáo (thường đã có sẵn):
%     \usepackage{amsmath, amssymb}   \usepackage{tikz}   \usepackage{graphicx}
% Dịch bằng XeLaTeX hoặc LuaLaTeX (có fontspec) để gõ tiếng Việt có dấu.
% ============================================================================

% ---- Kiểu vẽ dùng cho các hình trong chương (không cần thêm thư viện TikZ) ----
\tikzset{nut/.style={circle,draw,fill=white,minimum size=7mm,inner sep=1pt}}
% \tamgiac{tên}{x,y}{nhãn}: cây con vẽ bằng tam giác, (x,y) là đỉnh nối với cha
\newcommand{\tamgiac}[3]{%
  \coordinate (#1) at (#2);
  \draw[fill=gray!15] (#1) -- ++(-0.5,-0.9) -- ++(1,0) -- cycle;
  \node[font=\scriptsize] at ([yshift=-0.62cm]#1) {#3};}

\chapter{Thiết kế cấu trúc dữ liệu và giải thuật}

Chương này trình bày cơ sở lý thuyết và thiết kế của cây cân bằng AVL được chọn
làm cấu trúc dữ liệu chính cho hệ thống quản lý hàng tồn kho. Nội dung gồm lý do
chọn cấu trúc (mục~\ref{sec:avl-lydo}), định nghĩa và hệ số cân bằng
(mục~\ref{sec:avl-dinhnghia}), bốn phép quay (mục~\ref{sec:avl-quay}), các thao
tác thêm, xóa, tìm kiếm, duyệt và Range Search (mục~\ref{sec:avl-them}
đến mục~\ref{sec:avl-range}) và đánh giá độ phức tạp (mục~\ref{sec:avl-dophuctap}).

\section{Lý do chọn cây AVL}\label{sec:avl-lydo}

Danh mục mặt hàng của siêu thị thay đổi liên tục: mặt hàng mới được thêm vào,
mặt hàng ngừng kinh doanh bị xóa, số lượng tồn kho thay đổi sau mỗi lần nhập và
xuất. Đồng thời hệ thống phải tra cứu một mặt hàng theo mã gần như tức thì và
liệt kê các mặt hàng trong một khoảng mã. Các yêu cầu đó dẫn đến ba điều kiện
đối với cấu trúc dữ liệu: tìm, thêm và xóa nhanh; duyệt được theo thứ tự khóa;
và tra cứu được theo khoảng.

Cây nhị phân tìm kiếm (BST) thông thường đáp ứng điều kiện thứ hai và thứ ba,
nhưng độ phức tạp của nó phụ thuộc vào hình dạng cây. Nếu các mã hàng được nhập
theo thứ tự tăng dần, BST suy biến thành một danh sách liên kết có chiều cao
$n$ và mỗi thao tác tốn $O(n)$. Cây AVL giữ chiều cao luôn ở mức $O(\log n)$
bất kể thứ tự nhập, nên được chọn cho đề tài này.

\section{Định nghĩa cây AVL và hệ số cân bằng}\label{sec:avl-dinhnghia}

\subsection{Chiều cao và hệ số cân bằng}

Với mỗi nút $v$ của cây, ký hiệu $h(v)$ là chiều cao của cây con gốc $v$. Quy ước
dùng trong cài đặt: cây rỗng có chiều cao $0$, nút lá có chiều cao $1$, và
\[
  h(v) = 1 + \max\bigl(h(\text{con trái}),\, h(\text{con phải})\bigr).
\]
\textbf{Hệ số cân bằng} (balance factor) của nút $v$ là hiệu chiều cao hai cây con:
\[
  BF(v) = h(\text{con trái}) - h(\text{con phải}).
\]

\subsection{Định nghĩa}

\textbf{Cây AVL} là cây nhị phân tìm kiếm trong đó mọi nút $v$ đều thỏa
$BF(v) \in \{-1, 0, 1\}$. Nút có $BF > 1$ gọi là lệch trái, nút có $BF < -1$
gọi là lệch phải. Trong chương trình, mỗi nút \texttt{AVLNode} lưu sẵn trường
\texttt{height}, nên \texttt{getBalanceFactor} chỉ cần lấy hiệu hai trường này
mà không phải duyệt lại cây con.

\subsection{Chiều cao cây AVL}\label{sec:avl-chieucao}

Gọi $N(h)$ là số nút ít nhất của một cây AVL có chiều cao $h$. Cây ít nút nhất
có một cây con cao $h-1$ và cây con kia cao $h-2$ (lệch tối đa mà vẫn hợp lệ),
do đó
\[
  N(1) = 1,\quad N(2) = 2,\quad N(h) = N(h-1) + N(h-2) + 1 .
\]
Quy hồi này cùng dạng với dãy Fibonacci $F_1 = F_2 = 1,\ F_k = F_{k-1} + F_{k-2}$:
bằng quy nạp, $N(h) = F_{h+2} - 1$. Vì $F_{h+2} \ge \varphi^{h}$ với
$\varphi = \frac{1+\sqrt5}{2} \approx 1{,}618$, một cây AVL có $n$ nút và chiều
cao $h$ thỏa $n \ge \varphi^{h} - 1$, suy ra
\[
  h \;\le\; \log_{\varphi}(n+1) \;\approx\; 1{,}44 \,\log_2 (n+1).
\]
Như vậy chiều cao cây AVL luôn là $O(\log n)$, và mọi thao tác đi dọc một đường
từ gốc xuống lá đều có độ phức tạp $O(\log n)$. Bảng~\ref{tab:chieucao} cho
thấy giá trị này với các kích thước cây điển hình.

\begin{table}[htbp]
\centering
\begin{tabular}{|r|r|r|}
\hline
Số nút $n$ & Chiều cao tối đa của AVL ($\le 1{,}44\log_2(n+1)$) & Chiều cao tối đa của BST \\
\hline
$1\,000$   & $14$  & $1\,000$ \\
$10\,000$  & $19$  & $10\,000$ \\
$100\,000$ & $23$  & $100\,000$ \\
\hline
\end{tabular}
\caption{Giới hạn chiều cao của cây AVL so với trường hợp xấu nhất của BST.}
\label{tab:chieucao}
\end{table}

\section{Bốn phép quay cân bằng}\label{sec:avl-quay}

Sau khi thêm hoặc xóa một nút, một số nút trên đường đi từ nút đó lên gốc có thể
có $BF = \pm 2$. Gọi $z$ là nút lệch \emph{thấp nhất} (gần lá nhất), $y$ là con
của $z$ nằm ở phía cao hơn và $x$ là con của $y$ nằm ở phía cao hơn (hoặc bằng
nhau khi xóa). Tùy vị trí của $y$ và $x$ mà ta có bốn trường hợp. Trong các hình
dưới đây, $T_1, T_2, T_3, T_4$ là các cây con (vẽ bằng tam giác) và thứ tự khóa
từ trái sang phải được giữ nguyên trước và sau khi quay.

\subsection{Trường hợp LL: quay phải}

Nút $z$ lệch trái ($BF(z) = 2$) và con trái $y$ cũng lệch trái hoặc cân bằng
($BF(y) \ge 0$). Ta quay phải tại $z$: $y$ trở thành gốc của cây con, $z$ trở
thành con phải của $y$, cây con phải cũ của $y$ (tức $T_3$) được chuyển thành con
trái của $z$.

\begin{figure}[htbp]
\centering
\resizebox{\textwidth}{!}{%
\begin{tikzpicture}[every node/.style={font=\small}]
  % trước khi quay
  \node[nut] (z) at (3,0) {$z$};
  \node[nut] (y) at (1.5,-1.4) {$y$};
  \node[nut] (x) at (0,-2.8) {$x$};
  \tamgiac{a1}{-0.9,-4.0}{$T_1$}
  \tamgiac{a2}{0.9,-4.0}{$T_2$}
  \tamgiac{a3}{2.5,-2.8}{$T_3$}
  \tamgiac{a4}{4.5,-1.4}{$T_4$}
  \draw (z) -- (y);  \draw (z) -- (a4);
  \draw (y) -- (x);  \draw (y) -- (a3);
  \draw (x) -- (a1); \draw (x) -- (a2);
  \node[font=\scriptsize] at ([xshift=-9mm]z) {$BF=2$};
  % mũi tên
  \draw[->,thick] (5.8,-1.6) -- (7.2,-1.6)
      node[midway,above,font=\scriptsize,align=center] {quay phải\\tại $z$};
  % sau khi quay
  \node[nut] (y2) at (10,0) {$y$};
  \node[nut] (x2) at (8.5,-1.4) {$x$};
  \node[nut] (z2) at (11.5,-1.4) {$z$};
  \tamgiac{b1}{7.6,-2.8}{$T_1$}
  \tamgiac{b2}{9.4,-2.8}{$T_2$}
  \tamgiac{b3}{10.6,-2.8}{$T_3$}
  \tamgiac{b4}{12.4,-2.8}{$T_4$}
  \draw (y2) -- (x2); \draw (y2) -- (z2);
  \draw (x2) -- (b1); \draw (x2) -- (b2);
  \draw (z2) -- (b3); \draw (z2) -- (b4);
\end{tikzpicture}}
\caption{Trường hợp LL: quay phải tại $z$.}
\label{fig:quay-ll}
\end{figure}

\subsection{Trường hợp RR: quay trái}

Đối xứng với LL: $z$ lệch phải ($BF(z) = -2$) và con phải $y$ lệch phải hoặc cân
bằng ($BF(y) \le 0$). Ta quay trái tại $z$.

\begin{figure}[htbp]
\centering
\resizebox{\textwidth}{!}{%
\begin{tikzpicture}[every node/.style={font=\small}]
  \node[nut] (z) at (0,0) {$z$};
  \node[nut] (y) at (1.5,-1.4) {$y$};
  \node[nut] (x) at (3,-2.8) {$x$};
  \tamgiac{a1}{-1.5,-1.4}{$T_1$}
  \tamgiac{a2}{0.5,-2.8}{$T_2$}
  \tamgiac{a3}{2.1,-4.0}{$T_3$}
  \tamgiac{a4}{3.9,-4.0}{$T_4$}
  \draw (z) -- (a1); \draw (z) -- (y);
  \draw (y) -- (a2); \draw (y) -- (x);
  \draw (x) -- (a3); \draw (x) -- (a4);
  \node[font=\scriptsize] at ([xshift=9mm]z) {$BF=-2$};
  \draw[->,thick] (5.8,-1.6) -- (7.2,-1.6)
      node[midway,above,font=\scriptsize,align=center] {quay trái\\tại $z$};
  \node[nut] (y2) at (10,0) {$y$};
  \node[nut] (z2) at (8.5,-1.4) {$z$};
  \node[nut] (x2) at (11.5,-1.4) {$x$};
  \tamgiac{b1}{7.6,-2.8}{$T_1$}
  \tamgiac{b2}{9.4,-2.8}{$T_2$}
  \tamgiac{b3}{10.6,-2.8}{$T_3$}
  \tamgiac{b4}{12.4,-2.8}{$T_4$}
  \draw (y2) -- (z2); \draw (y2) -- (x2);
  \draw (z2) -- (b1); \draw (z2) -- (b2);
  \draw (x2) -- (b3); \draw (x2) -- (b4);
\end{tikzpicture}}
\caption{Trường hợp RR: quay trái tại $z$.}
\label{fig:quay-rr}
\end{figure}

\subsection{Trường hợp LR: quay trái rồi quay phải}

Nút $z$ lệch trái nhưng con trái $y$ lệch phải ($BF(y) < 0$), nên một phép quay
phải đơn tại $z$ chưa đủ. Ta quay trái tại $y$ để đưa về trường hợp LL, sau đó
quay phải tại $z$. Kết quả là $x$ (cháu) trở thành gốc cây con.

\begin{figure}[htbp]
\centering
\resizebox{\textwidth}{!}{%
\begin{tikzpicture}[every node/.style={font=\small}]
  \node[nut] (z) at (3,0) {$z$};
  \node[nut] (y) at (1.5,-1.4) {$y$};
  \node[nut] (x) at (2.6,-2.8) {$x$};
  \tamgiac{a1}{0.4,-2.8}{$T_1$}
  \tamgiac{a2}{1.7,-4.0}{$T_2$}
  \tamgiac{a3}{3.5,-4.0}{$T_3$}
  \tamgiac{a4}{4.5,-1.4}{$T_4$}
  \draw (z) -- (y);  \draw (z) -- (a4);
  \draw (y) -- (a1); \draw (y) -- (x);
  \draw (x) -- (a2); \draw (x) -- (a3);
  \node[font=\scriptsize] at ([xshift=-9mm]z) {$BF=2$};
  \draw[->,thick] (5.8,-1.6) -- (7.2,-1.6)
      node[midway,above,font=\scriptsize,align=center] {quay trái $y$,\\quay phải $z$};
  \node[nut] (x2) at (10,0) {$x$};
  \node[nut] (y2) at (8.5,-1.4) {$y$};
  \node[nut] (z2) at (11.5,-1.4) {$z$};
  \tamgiac{b1}{7.6,-2.8}{$T_1$}
  \tamgiac{b2}{9.4,-2.8}{$T_2$}
  \tamgiac{b3}{10.6,-2.8}{$T_3$}
  \tamgiac{b4}{12.4,-2.8}{$T_4$}
  \draw (x2) -- (y2); \draw (x2) -- (z2);
  \draw (y2) -- (b1); \draw (y2) -- (b2);
  \draw (z2) -- (b3); \draw (z2) -- (b4);
\end{tikzpicture}}
\caption{Trường hợp LR: quay trái tại $y$ rồi quay phải tại $z$.}
\label{fig:quay-lr}
\end{figure}

\subsection{Trường hợp RL: quay phải rồi quay trái}

Đối xứng với LR: $z$ lệch phải nhưng con phải $y$ lệch trái ($BF(y) > 0$). Ta quay
phải tại $y$ rồi quay trái tại $z$.

\begin{figure}[htbp]
\centering
\resizebox{\textwidth}{!}{%
\begin{tikzpicture}[every node/.style={font=\small}]
  \node[nut] (z) at (0,0) {$z$};
  \node[nut] (y) at (1.5,-1.4) {$y$};
  \node[nut] (x) at (0.4,-2.8) {$x$};
  \tamgiac{a1}{-1.5,-1.4}{$T_1$}
  \tamgiac{a2}{-0.4,-4.0}{$T_2$}
  \tamgiac{a3}{1.2,-4.0}{$T_3$}
  \tamgiac{a4}{2.6,-2.8}{$T_4$}
  \draw (z) -- (a1); \draw (z) -- (y);
  \draw (y) -- (x);  \draw (y) -- (a4);
  \draw (x) -- (a2); \draw (x) -- (a3);
  \node[font=\scriptsize] at ([xshift=9mm]z) {$BF=-2$};
  \draw[->,thick] (5.8,-1.6) -- (7.2,-1.6)
      node[midway,above,font=\scriptsize,align=center] {quay phải $y$,\\quay trái $z$};
  \node[nut] (x2) at (10,0) {$x$};
  \node[nut] (z2) at (8.5,-1.4) {$z$};
  \node[nut] (y2) at (11.5,-1.4) {$y$};
  \tamgiac{b1}{7.6,-2.8}{$T_1$}
  \tamgiac{b2}{9.4,-2.8}{$T_2$}
  \tamgiac{b3}{10.6,-2.8}{$T_3$}
  \tamgiac{b4}{12.4,-2.8}{$T_4$}
  \draw (x2) -- (z2); \draw (x2) -- (y2);
  \draw (z2) -- (b1); \draw (z2) -- (b2);
  \draw (y2) -- (b3); \draw (y2) -- (b4);
\end{tikzpicture}}
\caption{Trường hợp RL: quay phải tại $y$ rồi quay trái tại $z$.}
\label{fig:quay-rl}
\end{figure}

\subsection{Tổng hợp cách chọn phép quay}

Phép quay được chọn theo $BF$ của nút lệch $z$ và của con $y$ ở phía lệch
(bảng~\ref{tab:chonquay}). Trong chương trình, hàm \texttt{rotateRight} và
\texttt{rotateLeft} cài hai phép quay đơn; các trường hợp LR và RL được ghép từ
hai phép quay đơn. Sau mỗi phép quay, chiều cao của hai nút bị đổi chỗ được tính
lại theo thứ tự từ nút thấp lên nút cao.

\begin{table}[htbp]
\centering
\begin{tabular}{|c|c|l|l|}
\hline
$BF(z)$ & $BF(y)$ của con lệch & Trường hợp & Phép quay \\
\hline
$\ge 2$  & $\ge 0$ & LL & quay phải tại $z$ \\
$\ge 2$  & $<0$    & LR & quay trái tại $y$, quay phải tại $z$ \\
$\le -2$ & $\le 0$ & RR & quay trái tại $z$ \\
$\le -2$ & $>0$    & RL & quay phải tại $y$, quay trái tại $z$ \\
\hline
\end{tabular}
\caption{Chọn phép quay theo hệ số cân bằng.}
\label{tab:chonquay}
\end{table}

Mỗi phép quay chỉ đổi một số hữu hạn con trỏ nên tốn $O(1)$, và không làm đổi
thứ tự duyệt \emph{inorder} của các khóa: cây sau khi quay vẫn là cây nhị phân
tìm kiếm hợp lệ.

\section{Thêm nút và tái cân bằng}\label{sec:avl-them}

Thao tác thêm (\texttt{insert}) gồm hai pha, thực hiện trong cùng một hàm đệ quy
\texttt{insertNode}:
\begin{enumerate}
  \item \textbf{Đi xuống:} so sánh khóa mới với khóa của nút hiện tại, đi sang trái
        nếu nhỏ hơn và sang phải nếu lớn hơn, cho đến khi gặp vị trí rỗng thì tạo
        nút lá mới. Nếu gặp khóa trùng, thao tác dừng và trả về \texttt{false}, cây
        không thay đổi.
  \item \textbf{Đi lên:} trên đường quay lại gốc, mỗi nút cập nhật lại \texttt{height},
        tính $BF$, và nếu $|BF| = 2$ thì chọn một trong bốn phép quay ở
        bảng~\ref{tab:chonquay}.
\end{enumerate}

Sau khi thêm một nút, chiều cao của cây con bị lệch quay về đúng giá trị trước khi
thêm, nên \emph{nhiều nhất một lần tái cân bằng} (một phép quay đơn hoặc một cặp
phép quay kép) là đủ để khôi phục tính cân bằng cho toàn bộ cây. Vì đường đi
có độ dài $O(\log n)$ và mỗi phép quay tốn $O(1)$, thao tác thêm tốn
$O(\log n)$.

\section{Xóa nút và cân bằng lại}\label{sec:avl-xoa}

Thao tác xóa (\texttt{remove}) tìm nút theo khóa rồi xử lý ba trường hợp:
\begin{enumerate}
  \item \textbf{Nút lá:} xóa nút, con trỏ của cha trở thành rỗng.
  \item \textbf{Nút có một con:} xóa nút và nối con duy nhất của nó vào cha.
  \item \textbf{Nút có hai con:} tìm nút nhỏ nhất của cây con phải (nút kế tiếp theo
        thứ tự khóa, hàm \texttt{getMinValueNode}), chép khóa và dữ liệu của nó vào
        nút cần xóa, rồi xóa nút kế tiếp khỏi cây con phải. Nút kế tiếp không có
        con trái nên rơi vào trường hợp 1 hoặc 2.
\end{enumerate}
Trên đường quay lại gốc, mỗi nút cập nhật \texttt{height} và tái cân bằng giống
thao tác thêm. Có hai điểm khác so với thao tác thêm.
\begin{itemize}
  \item Sau khi xóa, con $y$ của nút lệch $z$ có thể có $BF(y) = 0$. Khi đó phải
        dùng phép quay \emph{đơn} (LL hoặc RR), vì quay kép sẽ không cân bằng lại
        cây. Điều kiện chọn quay kép chỉ là $BF(y)$ ngược chiều với $BF(z)$
        (bảng~\ref{tab:chonquay}).
  \item Một phép quay có thể làm chiều cao cây con giảm đi $1$, khiến nút tổ tiên
        phía trên tiếp tục bị lệch. Vì vậy xóa có thể cần tái cân bằng ở
        \emph{nhiều} nút trên đường đi, tối đa $O(\log n)$ phép quay.
\end{itemize}
Tổng độ phức tạp của thao tác xóa vẫn là $O(\log n)$: đi xuống tìm nút, tìm nút kế
tiếp và quay lại đều dọc theo một đường từ gốc xuống lá.

\section{Tìm kiếm và duyệt cây}\label{sec:avl-timkiem}

\subsection{Tìm kiếm theo khóa}

Hàm \texttt{search} dùng tính chất của cây nhị phân tìm kiếm: so sánh khóa cần tìm
với khóa của nút hiện tại, đi sang trái nếu nhỏ hơn, sang phải nếu lớn hơn, và
dừng khi bằng hoặc khi hết cây. Số bước không vượt quá chiều cao cây, tức
$O(\log n)$ nhờ tính chất ở mục~\ref{sec:avl-chieucao}. Hàm trả về con trỏ tới
dữ liệu thực tế của nút (hoặc \texttt{nullptr} nếu không có khóa), nên có thể cập
nhật trực tiếp, chẳng hạn số lượng tồn kho sau một lần nhập hoặc xuất, mà không
phải xóa rồi thêm lại.

\subsection{Duyệt trung thứ tự (LNR)}

Duyệt \emph{inorder} thăm cây con trái, rồi nút hiện tại, rồi cây con phải. Do tính
chất của cây nhị phân tìm kiếm, kết quả là danh sách dữ liệu theo khóa tăng dần.
Hàm \texttt{inorder} dùng thao tác này để lập danh sách mặt hàng theo mã tăng dần
cho các chức năng liệt kê và báo cáo, đồng thời là cách duy nhất cần dùng khi ghi
toàn bộ cây ra tệp theo thứ tự khóa. Mỗi nút được thăm đúng một lần nên độ phức
tạp là $O(n)$.

\section{Tra cứu theo phạm vi (Range Search)}\label{sec:avl-range}

Range Search trả về tất cả phần tử có khóa nằm trong đoạn $[lo, hi]$, theo thứ tự
tăng dần. Ví dụ: liệt kê các mặt hàng có mã từ \texttt{SP010} đến \texttt{SP050}.
Thay vì duyệt cả cây rồi lọc, thuật toán chỉ đi vào những nhánh có thể chứa kết
quả (hàm \texttt{rangeSearchNode}). Tại nút $v$ với khóa $k$:
\begin{itemize}
  \item nếu $lo < k$, cây con trái có thể chứa khóa trong đoạn, nên đi sang trái;
  \item nếu $lo \le k \le hi$, ghi nhận nút $v$ vào kết quả;
  \item nếu $k < hi$, cây con phải có thể chứa khóa trong đoạn, nên đi sang phải.
\end{itemize}
Thứ tự \emph{trái, nút, phải} bảo đảm kết quả đã được sắp tăng dần. Nhánh nào nằm
hoàn toàn ngoài đoạn bị bỏ qua ngay. Các nút được thăm chia thành hai loại: các
nút nằm trên hai đường từ gốc đến $lo$ và đến $hi$, có $O(\log n)$ nút, và các nút
nằm trong kết quả, có $k$ nút. Do đó thời gian là $O(\log n + k)$, thấp hơn nhiều
so với $O(n)$ khi duyệt toàn bộ nếu đoạn truy vấn nhỏ. Khi $lo > hi$ khoảng rỗng và
không có phần tử nào được trả về.

Cây được tổ chức theo \emph{khóa}, nên Range Search theo khoảng mã hàng dùng trực
tiếp khóa chính. Để tra cứu theo khoảng giá, nhiều mặt hàng có thể cùng giá trong
khi cây AVL không cho phép trùng khóa. Cách xử lý là dùng khóa ghép
$(\text{giá}, \text{mã hàng})$ với phép so sánh giá trước, mã hàng sau, và dùng cận
dưới $(g_{\min}, \text{rỗng})$ và cận trên $(g_{\max}, \text{ký tự lớn nhất})$
cho $lo$ và $hi$.

\section{Đánh giá độ phức tạp}\label{sec:avl-dophuctap}

Bảng~\ref{tab:phuctap} tổng hợp độ phức tạp thời gian của các thao tác, với $n$ là
số nút của cây và $k$ là số kết quả của Range Search. Mọi giá trị đều suy ra từ
chiều cao $h = O(\log n)$ ở mục~\ref{sec:avl-chieucao}.

\begin{table}[htbp]
\centering
\begin{tabular}{|l|c|c|l|}
\hline
Thao tác & Trung bình & Xấu nhất & Ghi chú \\
\hline
Thêm (\texttt{insert})       & $O(\log n)$     & $O(\log n)$     & tối đa 1 lần tái cân bằng \\
Xóa (\texttt{remove})        & $O(\log n)$     & $O(\log n)$     & có thể quay ở nhiều nút \\
Tìm (\texttt{search})        & $O(\log n)$     & $O(\log n)$     & \\
Duyệt (\texttt{inorder})     & $O(n)$          & $O(n)$          & mỗi nút thăm một lần \\
Range Search                 & $O(\log n + k)$ & $O(\log n + k)$ & $k$ là số kết quả \\
Bộ nhớ                       & $O(n)$          & $O(n)$          & thêm 1 trường \texttt{height} mỗi nút \\
\hline
\end{tabular}
\caption{Độ phức tạp các thao tác trên cây AVL.}
\label{tab:phuctap}
\end{table}

Với BST thông thường, các thao tác thêm, xóa và tìm có độ phức tạp trung bình
$O(\log n)$ nhưng xấu nhất là $O(n)$ khi dữ liệu nhập theo thứ tự. Cây AVL đánh đổi
một chút chi phí cho việc quay và một trường \texttt{height} trên mỗi nút để có
bảo đảm $O(\log n)$ trong \emph{mọi} trường hợp, điều cần thiết khi mã hàng thường
được cấp theo thứ tự tăng dần.

Kết quả đo thực tế trên bộ test (chương~4) xác nhận các giá trị này: với $10^3$,
$10^4$ và $10^5$ khóa ngẫu nhiên, chiều cao đo được lần lượt là $12$, $16$ và $20$,
đều thấp hơn giới hạn lý thuyết $1{,}44\log_2(n+1)$ ở bảng~\ref{tab:chieucao}.
